\documentclass[10pt,a4paper]{amsart}
\usepackage{amsmath,amssymb,amsthm,graphicx}
% Retained minimal preamble for the Stacks Project chapter flat.tex (More on Flatness).
% Verbatim excerpt of the chapter preamble: Xy-pic setup (native preamble lines 14-18) and the native
% theorem declarations with the mathematical macros (native preamble lines 154-217). Copied mechanically;
% nothing added, removed or redefined. The wrapper class and the amsmath/amssymb/amsthm/graphicx base
% packages are supplied by the assembler, so the class/externaldocument part of the native preamble is omitted.
% For commutative diagrams we use Xy-pic
\usepackage[all]{xy}

% We use 2cell for 2-commutative diagrams.
\xyoption{2cell}
% Theorem environments.
%
\theoremstyle{plain}
\newtheorem{theorem}[subsection]{Theorem}
\newtheorem{proposition}[subsection]{Proposition}
\newtheorem{lemma}[subsection]{Lemma}

\theoremstyle{definition}
\newtheorem{definition}[subsection]{Definition}
\newtheorem{example}[subsection]{Example}
\newtheorem{exercise}[subsection]{Exercise}
\newtheorem{situation}[subsection]{Situation}

\theoremstyle{remark}
\newtheorem{remark}[subsection]{Remark}
\newtheorem{remarks}[subsection]{Remarks}

\numberwithin{equation}{subsection}

% Macros
%
\def\lim{\mathop{\mathrm{lim}}\nolimits}
\def\colim{\mathop{\mathrm{colim}}\nolimits}
\def\Spec{\mathop{\mathrm{Spec}}}
\def\Hom{\mathop{\mathrm{Hom}}\nolimits}
\def\Ext{\mathop{\mathrm{Ext}}\nolimits}
\def\SheafHom{\mathop{\mathcal{H}\!\mathit{om}}\nolimits}
\def\SheafExt{\mathop{\mathcal{E}\!\mathit{xt}}\nolimits}
\def\Sch{\mathit{Sch}}
\def\Mor{\mathop{\mathrm{Mor}}\nolimits}
\def\Ob{\mathop{\mathrm{Ob}}\nolimits}
\def\Sh{\mathop{\mathit{Sh}}\nolimits}
\def\NL{\mathop{N\!L}\nolimits}
\def\CH{\mathop{\mathrm{CH}}\nolimits}
\def\proetale{{pro\text{-}\acute{e}tale}}
\def\etale{{\acute{e}tale}}
\def\QCoh{\mathit{QCoh}}
\def\Ker{\mathop{\mathrm{Ker}}}
\def\Im{\mathop{\mathrm{Im}}}
\def\Coker{\mathop{\mathrm{Coker}}}
\def\Coim{\mathop{\mathrm{Coim}}}

% Boxtimes
%
\DeclareMathSymbol{\boxtimes}{\mathbin}{AMSa}{"02}

%
% Macros for moduli stacks/spaces
%
\def\QCohstack{\mathcal{QC}\!\mathit{oh}}
\def\Cohstack{\mathcal{C}\!\mathit{oh}}
\def\Spacesstack{\mathcal{S}\!\mathit{paces}}
\def\Quotfunctor{\mathrm{Quot}}
\def\Hilbfunctor{\mathrm{Hilb}}
\def\Curvesstack{\mathcal{C}\!\mathit{urves}}
\def\Polarizedstack{\mathcal{P}\!\mathit{olarized}}
\def\Complexesstack{\mathcal{C}\!\mathit{omplexes}}
% \Pic is the operator that assigns to X its picard group, usage \Pic(X)
% \Picardstack_{X/B} denotes the Picard stack of X over B
% \Picardfunctor_{X/B} denotes the Picard functor of X over B
\def\Pic{\mathop{\mathrm{Pic}}\nolimits}
\def\Picardstack{\mathcal{P}\!\mathit{ic}}
\def\Picardfunctor{\mathrm{Pic}}
\def\Deformationcategory{\mathcal{D}\!\mathit{ef}}

% --- appended 2026-10-12 (Chief): the project own definitions of the reference/slogan/history
% display environments, copied mechanically from the shared Stacks preamble lines 8-12
% (\usepackage{verbatim} plus the three \newenvironment lines, which the chapter preamble keeps
% verbatim from the project preamble). Nothing added, removed or redefined: these environments
% carry no mathematical content and the project itself typesets them as comments. Without these
% three lines a counted statement whose own text uses them fails with Environment slogan undefined.
\usepackage{verbatim}
\newenvironment{reference}{\comment}{\endcomment}
\newenvironment{slogan}{\comment}{\endcomment}
\newenvironment{history}{\comment}{\endcomment}
\usepackage{cleveref}
\crefname{originaltheorem}{original theorem}{original theorems}
\Crefname{originaltheorem}{Original theorem}{Original theorems}
\makeatletter
\newlabel{algebra-theorem-projective-free-over-local-ring}{{0593}{source}{}{sourcepointer-1}{}}
\newlabel{algebra-theorem-projective-free-over-local-ring@cref}{{[originaltheorem][0][]0593}{[1][0][]source}{}{}{}}
\makeatother
\crefname{originallemma}{original lemma}{original lemmas}
\Crefname{originallemma}{Original lemma}{Original lemmas}
\makeatletter
\newlabel{lemma-flat-finite-presentation-affine-neighbourhood-projective}{{05ME}{source}{}{sourcepointer-2}{}}
\newlabel{lemma-flat-finite-presentation-affine-neighbourhood-projective@cref}{{[originallemma][0][]05ME}{[1][0][]source}{}{}{}}
\makeatother
\crefname{originalproposition}{original proposition}{original propositions}
\Crefname{originalproposition}{Original proposition}{Original propositions}
\makeatletter
\newlabel{proposition-finite-type-flat-at-point}{{05I5}{source}{}{sourcepointer-3}{}}
\newlabel{proposition-finite-type-flat-at-point@cref}{{[originalproposition][0][]05I5}{[1][0][]source}{}{}{}}
\makeatother
\title{Free global sections on an affine etale neighbourhood for a finite type module flat at a point}
\author{The Stacks Project Authors}
\date{stacks-project:flat.tex:local-snapshot-20261011; GNU FDL 1.2 or later}
\begin{document}
\maketitle
\noindent\textit{Editorial note.} All delivered mathematical text is the author's own, copied from the pinned Stacks Project chapter file flat.tex; no statement or proof step was rewritten, reordered or removed.\par\medskip
\noindent\textit{Editorial note.} Preamble provenance: the chapter's own retained minimal preamble (sources/stacks-flat/retained-preamble.tex, lines 1-85) is a verbatim excerpt of the native Xy-pic setup and of the native theorem declarations and mathematical macros, followed by the project's own reference/slogan/history display-environment definitions copied mechanically from the shared preamble; the wrapper class and the amsmath/amssymb/amsthm/graphicx base packages are supplied by the assembler, and extra\_packages is left empty because everything needed is already in the retained preamble. No macro was invented or redefined.\par\medskip
\noindent\textit{Editorial note.} Same-chapter pointers are located at their \textbackslash{}label lines inside the pinned chapter file; cross-chapter pointers are located in the project's own tag list staged in the snapshot (sources/stacks-chow/tags.txt), at the verbatim lines carrying their canonical tags.\par\medskip
\noindent\textit{Allowed standard prerequisites (editorial).}\begin{itemize}
\item The affine neighbourhood statement in the finite presentation case, and freeness of a projective module over a local ring. Same chapter flat.tex, label lemma-flat-finite-presentation-affine-neighbourhood-projective (proof line 5752, claimed as S14-algebra-projective-neighbourhood); Algebra, label algebra-theorem-projective-free-over-local-ring (proof line 5754). The first provides the affine etale neighbourhood with projective global sections once the module is of finite presentation, and the second turns projectivity over the local base into freeness.
\item The local structure statement giving a finitely presented pullback after an etale neighbourhood. Same chapter flat.tex, label proposition-finite-type-flat-at-point, invoked at proof line 5750. Allows the reduction to a module of finite presentation by passing to an affine neighbourhood on the base and an open of the base change.
\end{itemize}\medskip
\ifdefined\hypertarget\hypertarget{sourcepointer-1}{}\fi
\noindent\textit{Source reference.} Original theorem 0593: sources/stacks-chow/tags.txt, source lines 6801--6801. Navigation pointer for \textbackslash{}ref\{algebra-theorem-projective-free-over-local-ring\} used inside the counted span (line 5754); the argument that consumes it is retained in the span.\par\medskip
\ifdefined\hypertarget\hypertarget{sourcepointer-2}{}\fi
\noindent\textit{Source reference.} Original lemma 05ME: sources/stacks-flat/pinned-original-source.tex, source lines 5662--5662. Navigation pointer for \textbackslash{}ref\{lemma-flat-finite-presentation-affine-neighbourhood-projective\} used inside the counted span (line 5752); the argument that consumes it is retained in the span.\par\medskip
\ifdefined\hypertarget\hypertarget{sourcepointer-3}{}\fi
\noindent\textit{Source reference.} Original proposition 05I5: sources/stacks-flat/pinned-original-source.tex, source lines 2477--2477. Navigation pointer for \textbackslash{}ref\{proposition-finite-type-flat-at-point\} used inside the counted span (line 5750); the argument that consumes it is retained in the span.\par\medskip
\begin{lemma}
\label{lemma-flat-finite-type-affine-neighbourhood-projective}
Let $f : X \to S$ be a morphism which is locally of finite type.
Let $\mathcal{F}$ be a quasi-coherent $\mathcal{O}_X$-module which is
of finite type. Let $x \in X$ with $s = f(x) \in S$.
If $\mathcal{F}$ is flat at $x$ over $S$ there exists an affine
elementary \'etale neighbourhood $(S', s') \to (S, s)$ and
an affine open $U' \subset X \times_S \Spec(\mathcal{O}_{S', s'})$
which contains $x' = (x, s')$ such that
$\Gamma(U', \mathcal{F}|_{U'})$ is a free
$\mathcal{O}_{S', s'}$-module.
\end{lemma}

\begin{proof}
The question is Zariski local on $X$ and $S$. Hence we may assume
that $X$ and $S$ are affine. Then we can find a closed immersion
$i : X \to \mathbf{A}^n_S$ over $S$. It is clear that it suffices to
prove the lemma for the sheaf $i_*\mathcal{F}$ on $\mathbf{A}^n_S$
and the point $i(x)$. In this way we reduce to the case where $X \to S$ is
of finite presentation. After replacing $S$ by
$\Spec(\mathcal{O}_{S', s'})$ and $X$ by an open of
$X \times_S \Spec(\mathcal{O}_{S', s'})$ we may assume that
$\mathcal{F}$ is of finite presentation, see
Proposition \ref{proposition-finite-type-flat-at-point}.
In this case we may appeal to
Lemma \ref{lemma-flat-finite-presentation-affine-neighbourhood-projective}
and
Algebra, Theorem \ref{algebra-theorem-projective-free-over-local-ring}
to conclude.
\end{proof}

\end{document}
